\subsection{Рекомендації щодо трансформації навчальних дисциплін}\label{subsec:3.2}

Структурно-функціональна модель (п.~\ref{subsec:3.1}) визначає загальну архітектуру формування ІКК. Цей підрозділ конкретизує змістовий та організаційно-технологічний блоки моделі через розробку рекомендацій щодо трансформації навчальних дисциплін для кожного типу курсів, а також представляє інтегровану модель 240-кредитного навчального плану з наскрізною ШІ-інтеграцією.

Рекомендації розроблено на основі: (а)~аналізу 132 силабусів КДПУ (п.~\ref{subsec:2.2}), який виявив нульове покриття ШІ-тематики; (б)~міжнародного бенчмаркінгу 7 провідних програм (п.~\ref{subsec:2.3}); (в)~навчального плану, запропонованого на основі аналізу прогалин; (г)~оглядового дослідження 156 наукових праць~\cite{Musiienko2026ScopingReview}, що визначило ключові тематичні кластери інтеграції ШІ у дизайн-освіту.

\subsubsection{Принципи трансформації навчальних дисциплін}\label{sssec:3.2.1}

Перед тим як перейти до конкретних рекомендацій для кожного типу дисциплін, визначимо загальні принципи трансформації, що випливають зі структурно-функціональної моделі:

\begin{enumerate}
    \item \textbf{Не створення нових дисциплін, а оновлення існуючих}~--- як основний механізм. З огляду на ригідність українських навчальних планів (фактична відсутність вибіркових компонентів, виявлена у п.~\ref{subsec:2.1}), основним шляхом інтеграції ШІ є оновлення змісту чинних дисциплін, а не додавання нових. Водночас модель передбачає введення 2--3 нових ШІ-орієнтованих дисциплін на місце вибіркових компонентів або через перерозподіл кредитів.

    \item \textbf{Диференціація за рівнем ШІ-інтеграції}. Для кожної дисципліни визначається один із трьох рівнів інтеграції:
    \begin{itemize}
        \item \textit{базовий}~--- ознайомлення з можливостями ШІ у контексті дисципліни (1--2 заняття);
        \item \textit{інструментальний}~--- систематичне використання ШІ-інструментів як частини робочого процесу (модуль 4--8 годин);
        \item \textit{трансформаційний}~--- ШІ як центральний елемент дисципліни, що змінює її зміст та методику (наскрізна інтеграція).
    \end{itemize}

    \item \textbf{Збереження фундаменту}. Традиційні навички (рисунок, живопис, композиція, макетування) зберігаються як основа дизайнерської підготовки. ШІ \textit{доповнює}, а не замінює ці навички, що підтверджується міжнародною практикою (п.~\ref{subsec:2.3}).

    \item \textbf{Покомпонентна інтеграція ІКК}. Кожна рекомендація спрямована на формування конкретних компонентів ІКК (інформаційно-аналітичного, комунікативного, технологічного, рефлексивного) або їх комбінації.
\end{enumerate}

\subsubsection{Трансформація студійних курсів}\label{sssec:3.2.2}

Студійні курси (Живопис, Рисунок, Композиція, Колористика) є ядром дизайнерської підготовки та відповідають <<підписовій педагогіці>> студійної моделі~\cite{Musiienko2026ScopingReview}. Вони формують фундаментальні навички візуального мислення, що не можуть бути замінені ШІ-інструментами. Водночас вони є потужним середовищем для формування рефлексивного та інформаційно-аналітичного компонентів ІКК через модифіковану критику (crit).

\begin{table}[htbp]
\centering
\caption{Трансформація студійних курсів для формування ІКК}\label{tab:studio-transform}
\small
\begin{tabularx}{\textwidth}{@{}p{0.14\textwidth}p{0.25\textwidth}p{0.30\textwidth}p{0.23\textwidth}@{}}
\toprule
\textbf{Дисципліна} & \textbf{Поточний стан} & \textbf{Рекомендовані зміни} & \textbf{Компоненти ІКК} \\
\midrule
Рисунок і живопис & Традиційна техніка; олівець, акварель, олія; натюрморт, пленер & Базовий рівень ШІ: порівняння ручних робіт з ШІ-генерацією (Midjourney <<in the style of...>>); рефлексія щодо автентичності; ШІ як інструмент пошуку референсів & Рефлексивний, Інформаційно-аналітичний \\
\addlinespace
Основи композиції & Формальна композиція; ритм, контраст, домінанта; ручне виконання & Інструментальний рівень: генерація 20+ варіантів композиції через ШІ для порівняльного аналізу; вправи <<ручне vs. ШІ>>; документація процесу & Технологічний, Інформаційно-аналітичний, Рефлексивний \\
\addlinespace
Колористика & Теорія кольору; кольорові гармонії; ручне змішування & Базовий рівень: ШІ-генерація кольорових палітр (Adobe Color AI, Coolors AI); аналіз кольорових рішень у ШІ-зображеннях; етика кольорового bias & Технологічний, Рефлексивний \\
\bottomrule
\end{tabularx}
\end{table}

Ключовим методичним прийомом для студійних курсів є \textit{порівняльний аналіз <<ручне vs. ШІ>>}: студент виконує завдання традиційним способом, потім генерує аналогічне за допомогою ШІ-інструменту, після чого порівнює результати та рефлексує щодо сильних і слабких сторін кожного підходу. Цей прийом одночасно формує технологічний (робота з ШІ), інформаційно-аналітичний (порівняльний аналіз) та рефлексивний (обґрунтування) компоненти ІКК, не порушуючи основну мету курсу~--- формування ручних навичок.

\textbf{Приклад завдання для Композиції.} Тема: <<Ритмічна композиція>>. Етап~1: студент створює ручну композицію з ритмічними елементами (2 години). Етап~2: студент формулює промпти для Midjourney для створення аналогічної за структурою композиції, ітеративно вдосконалюючи промпт (1 година). Етап~3: студент порівнює результати, аналізуючи: точність передачі задуму, виразність, оригінальність, контроль над результатом. Етап~4: публічна критика з обговоренням обох підходів. Рефлексивний запис у щоденнику. Формовані компоненти ІКК: технологічний (промпт-інженерія), інформаційно-аналітичний (порівняльний аналіз), рефлексивний (обґрунтування), комунікативний (критика).

\subsubsection{Трансформація курсів цифрового дизайну}\label{sssec:3.2.3}

Курси цифрового дизайну (Комп'ютерна графіка/Adobe Suite, 3D-моделювання, Figma та прототипування, Моушн-дизайн) є природним середовищем для глибокої інтеграції ШІ, оскільки вони вже працюють у цифровому середовищі. Для цих курсів рекомендовано інструментальний або трансформаційний рівень інтеграції.

\begin{table}[htbp]
\centering
\caption{Трансформація курсів цифрового дизайну}\label{tab:digital-transform}
\small
\begin{tabularx}{\textwidth}{@{}p{0.14\textwidth}p{0.24\textwidth}p{0.31\textwidth}p{0.23\textwidth}@{}}
\toprule
\textbf{Дисципліна} & \textbf{Поточний стан} & \textbf{Рекомендовані зміни} & \textbf{Компоненти ІКК} \\
\midrule
Комп'ютерна графіка (Adobe Suite) & Photoshop, Illustrator, InDesign; ручні операції обробки & Інструментальний: Adobe Firefly генерація та редагування; Generative Fill/Expand; ШІ-ретуш; автоматизація рутинних операцій через ШІ; модуль промпт-інженерії & Технологічний, Інформаційно-аналітичний \\
\addlinespace
3D-моделювання & Blender/3ds Max; полігональне моделювання; текстурування & Інструментальний: ШІ-генерація 3D (Meshy, Luma AI); ШІ-текстурування; порівняння ручного та ШІ-моделювання; обговорення обмежень & Технологічний, Рефлексивний \\
\addlinespace
Figma та прототипування & Figma; компоненти; прототипи; варіанти & Інструментальний: Figma AI для генерації макетів; ШІ-плагіни (Magician, Relume); порівняння ШІ-прототипів з ручними; UX-дослідження з ШІ & Технологічний, Комунікативний \\
\addlinespace
Моушн-дизайн & After Effects; анімація; motion graphics & Трансформаційний: RunwayML для генеративного відео; ШІ-анімація; промпт-режисура; етика deepfake у дизайні; порівняння підходів & Технологічний, Рефлексивний, Комунікативний \\
\bottomrule
\end{tabularx}
\end{table}

Для курсів цифрового дизайну ключовим методичним підходом є \textit{ШІ-аугментований робочий процес}: студент опановує класичний інструмент (Photoshop, Blender, Figma), а потім вивчає, як ШІ-функції інтегруються у цей інструмент для прискорення та покращення роботи. Це забезпечує формування технологічного компоненту ІКК не як ізольованого навику, а як інтегральної частини професійного робочого процесу.

\textbf{Приклад модуля для Комп'ютерної графіки (Adobe Suite).} Тема: <<ШІ-інструменти в Adobe Photoshop>>. Заняття~1 (2 год.): Generative Fill~--- генеративне заповнення як інструмент фотомонтажу. Практика: розширення зображень, видалення об'єктів, додавання елементів через текстові промпти. Заняття~2 (2 год.): Adobe Firefly~--- генерація зображень з тексту та стилістичні трансформації. Практика: створення візуальних концепцій для проєкту. Заняття~3 (2 год.): Інтеграція ШІ у професійний робочий процес~--- коли використовувати ШІ, а коли~--- традиційні інструменти. Порівняльний аналіз результатів. Заняття~4 (2 год.): Етичні аспекти~--- авторське право на ШІ-генерований контент, маніпулятивне використання, deepfake. Рефлексивна дискусія. Загальний обсяг: 8 годин (4 заняття), інструментальний рівень інтеграції.

\subsubsection{Трансформація теоретичних курсів}\label{sssec:3.2.4}

Теоретичні курси (Історія мистецтва та дизайну, Дизайн-мислення, Естетика) формують концептуальну базу дизайнерської діяльності. Для них рекомендовано базовий або інструментальний рівень ШІ-інтеграції з акцентом на інформаційно-аналітичний та рефлексивний компоненти ІКК.

\begin{table}[htbp]
\centering
\caption{Трансформація теоретичних курсів}\label{tab:theory-transform}
\small
\begin{tabularx}{\textwidth}{@{}p{0.14\textwidth}p{0.24\textwidth}p{0.31\textwidth}p{0.23\textwidth}@{}}
\toprule
\textbf{Дисципліна} & \textbf{Поточний стан} & \textbf{Рекомендовані зміни} & \textbf{Компоненти ІКК} \\
\midrule
Історія мистецтва та дизайну & Хронологічний огляд; стилі та напрямки; візуальний аналіз & Базовий: ШІ-генерація <<у стилі>> історичних напрямків; критичний аналіз точності ШІ-інтерпретацій; дискусія <<Чи може ШІ створити мистецтво?>>; візуальний пошук через ШІ & Інформаційно-аналітичний, Рефлексивний \\
\addlinespace
Дизайн-мислення & Design Thinking процес; емпатія, ідеація, прототипування & Інструментальний: ШІ на етапі ідеації (генерація варіантів); ШІ для швидкого прототипування; аналіз впливу ШІ на Design Thinking процес; етика ШІ у дизайні для людей & Інформаційно-аналітичний, Технологічний, Рефлексивний \\
\addlinespace
Дослідження в дизайні & Методологія дизайн-досліджень; якісні та кількісні методи & Інструментальний: ШІ для аналізу даних UX-досліджень; ChatGPT для систематизації якісних даних; критична оцінка ШІ-аналітики; етика ШІ у дослідженнях & Інформаційно-аналітичний, Рефлексивний \\
\bottomrule
\end{tabularx}
\end{table}

Для теоретичних курсів ключовим є метод \textit{критичного аналізу ШІ-контенту}: студенти генерують ШІ-контент, пов'язаний з темою курсу (наприклад, зображення <<у стилі Баухаусу>> для Історії мистецтва), а потім критично аналізують точність, стереотипність та обмеження ШІ-інтерпретації. Це формує одночасно інформаційно-аналітичний та рефлексивний компоненти ІКК.

\subsubsection{Трансформація професійних курсів}\label{sssec:3.2.5}

Професійні курси (UX/UI дизайн, Типографіка, Графічний дизайн, Ілюстрація, Дизайн бренду, Інформаційний дизайн) є ключовим полем для формування технологічного компоненту ІКК, оскільки вони безпосередньо моделюють професійну діяльність дизайнера. Для них рекомендовано інструментальний або трансформаційний рівень інтеграції.

\begin{table}[htbp]
\centering
\caption{Трансформація професійних курсів}\label{tab:professional-transform}
\small
\begin{tabularx}{\textwidth}{@{}p{0.14\textwidth}p{0.24\textwidth}p{0.31\textwidth}p{0.23\textwidth}@{}}
\toprule
\textbf{Дисципліна} & \textbf{Поточний стан} & \textbf{Рекомендовані зміни} & \textbf{Компоненти ІКК} \\
\midrule
UX/UI: Основи & Wireframing; user flows; usability & Інструментальний: ШІ для генерації wireframes (Figma AI, Uizard); ШІ-аналіз user flows; A/B тестування з ШІ; промпт-інженерія для UI-копірайтингу & Технологічний, Інформаційно-аналітичний \\
\addlinespace
UX/UI: Продвинутий & Складні інтерфейси; дослідницькі методи; дизайн-системи & Трансформаційний: повний UX/UI процес з ШІ~--- від дослідження до прототипу; ШІ в UX-аналітиці; дизайн для ШІ-інтерфейсів (conversational UI) & Технологічний, Комунікативний, Рефлексивний \\
\addlinespace
Типографіка & Шрифтознавство; верстка; типографічна ієрархія & Базовий: ШІ-генерація шрифтових пар; аналіз типографіки в ШІ-зображеннях; можливості та обмеження ШІ у шрифтовому дизайні & Технологічний, Інформаційно-аналітичний \\
\addlinespace
Графічний дизайн & Фірмовий стиль; плакат; поліграфія & Інструментальний: ШІ для генерації логотипів та варіантів; промпт-інженерія для стилістичної консистентності; ШІ у робочому процесі дизайн-студії & Технологічний, Комунікативний \\
\addlinespace
Ілюстрація & Авторська ілюстрація; техніки; стилізація & Інструментальний: ШІ-ілюстрація (Midjourney, DALL-E); порівняння авторського та ШІ-стилю; етика авторства; ШІ як інструмент пошуку ідей & Технологічний, Рефлексивний \\
\addlinespace
Дизайн бренду & Брендинг; позиціонування; візуальна ідентичність & Інструментальний: ШІ для генерації концепцій бренду; ШІ-аналіз конкурентів; ChatGPT для генерації brand voice; мудборди через ШІ & Технологічний, Інформаційно-аналітичний, Комунікативний \\
\addlinespace
Інформаційний дизайн & Інфографіка; візуалізація даних; діаграми & Інструментальний: ШІ для візуалізації даних; ChatGPT для аналізу даних; автоматизація інфографіки; критичний аналіз ШІ-візуалізацій & Технологічний, Інформаційно-аналітичний \\
\bottomrule
\end{tabularx}
\end{table}

Для професійних курсів ключовим є метод \textit{ШІ-інтегрованого проєктного навчання}: кожен проєкт включає обов'язковий етап використання ШІ-інструментів з документацією процесу (промпти, ітерації, обґрунтування вибору). Студент подає не лише фінальний результат, а й <<ШІ-щоденник проєкту>>, що документує взаємодію з ШІ. Промпт-інженерія тут виступає як нова форма професійної комунікації <<людина--ШІ>>, що відповідає виявленому у скопінг-огляді~\cite{Musiienko2026ScopingReview} зсуву до ролі <<дизайнера-куратора>>.

\textbf{Приклад проєктного завдання для Графічного дизайну.} Тема: <<Розробка фірмового стилю>>. Бриф: розробити фірмовий стиль для фіктивного стартапу. Етапи: 1)~Дослідження: пошук аналогів, аналіз тенденцій (ШІ: ChatGPT для аналізу ринку, Midjourney для генерації мудбордів). 2)~Ідеація: генерація 50+ варіантів логотипу (ручні ескізи + Midjourney/DALL-E генерація). 3)~Відбір та доопрацювання: відбір 3--5 найкращих варіантів, доопрацювання у Illustrator з використанням Adobe Firefly. 4)~Система: розробка guidelines (Figma AI для генерації макетів носіїв). 5)~Презентація: публічна критика з демонстрацією всього процесу, включаючи ШІ-компонент. 6)~Рефлексія: аналіз ролі ШІ у проєкті~--- де ШІ допоміг, де завадив, етичні питання. Формовані компоненти ІКК: усі чотири.

\subsubsection{Трансформація практики та стажування}\label{sssec:3.2.6}

Навчальна та виробнича практики є критичним елементом формування ІКК, оскільки вони забезпечують контакт з реальним професійним середовищем. Аналіз опитувань роботодавців (п.~\ref{subsec:2.2}) підтвердив, що саме роботодавці є найактивнішими ініціаторами ШІ-інтеграції, прямо рекомендуючи <<використання штучного інтелекту у сфері дизайну>>.

\begin{table}[htbp]
\centering
\caption{Трансформація практик та стажувань}\label{tab:practice-transform}
\small
\begin{tabularx}{\textwidth}{@{}p{0.14\textwidth}p{0.24\textwidth}p{0.31\textwidth}p{0.23\textwidth}@{}}
\toprule
\textbf{Вид практики} & \textbf{Поточний стан} & \textbf{Рекомендовані зміни} & \textbf{Компоненти ІКК} \\
\midrule
Навчальна практика (2 курс) & Ознайомча; спостереження за роботою дизайнерів & Інструментальний: обов'язкова фіксація використання ШІ на підприємстві; виконання ШІ-проєкту під керівництвом наставника; порівняльний звіт & Технологічний, Інформаційно-аналітичний \\
\addlinespace
Виробнича практика (3 курс) & Самостійна робота на підприємстві; виконання завдань & Трансформаційний: вибір баз практики з урахуванням ШІ-використання; самостійний ШІ-проєкт; документація ШІ-робочих процесів підприємства; ШІ-портфоліо & Усі компоненти \\
\addlinespace
Переддипломна практика (4 курс) & Збір матеріалу для дипломного проєкту & Трансформаційний: ШІ як дослідницький та проєктний інструмент; документація ШІ-стратегії для дипломного проєкту; рефлексивний звіт про роль ШІ & Усі компоненти \\
\bottomrule
\end{tabularx}
\end{table}

Для практик ключовим є формування \textit{переліку баз практики, що використовують ШІ-інструменти}. У зв'язку з тим, що роботодавці програми <<Дизайн одягу>> прямо рекомендували інтеграцію ШІ (п.~\ref{subsec:2.2}), КДПУ може ініціювати партнерства з дизайн-студіями, що активно впроваджують ШІ у свої робочі процеси.

Вимоги до звіту з виробничої практики мають бути доповнені: окрім традиційних розділів (опис підприємства, виконані завдання, висновки), звіт має включати: (а)~опис ШІ-інструментів, які використовуються на підприємстві; (б)~аналіз ефективності ШІ у робочих процесах дизайну; (в)~рефлексію щодо ШІ-готовності підприємства та власної ШІ-компетентності; (г)~приклади власних робіт, виконаних з використанням ШІ.

\subsubsection{Трансформація загальноосвітніх курсів}\label{sssec:3.2.7}

Загальноосвітні курси (Іноземна мова, Філософія, Академічне письмо, Дизайн та підприємництво) мають опосередкований, але важливий вплив на формування ІКК, зокрема комунікативного та рефлексивного компонентів.

\begin{table}[htbp]
\centering
\caption{Трансформація загальноосвітніх курсів}\label{tab:general-transform}
\small
\begin{tabularx}{\textwidth}{@{}p{0.14\textwidth}p{0.24\textwidth}p{0.31\textwidth}p{0.23\textwidth}@{}}
\toprule
\textbf{Дисципліна} & \textbf{Поточний стан} & \textbf{Рекомендовані зміни} & \textbf{Компоненти ІКК} \\
\midrule
Іноземна мова & Загальна англійська; граматика та лексика & Інструментальний: English for Design~--- фахова англійська з акцентом на ШІ-термінологію; робота з англомовною документацією ШІ-інструментів; промпт-інженерія англійською; ChatGPT як мовний інструмент & Комунікативний, Технологічний \\
\addlinespace
Академічне письмо & Основи наукового стилю; структура тексту & Базовий: ШІ як асистент написання (ChatGPT, Grammarly AI); критичний аналіз ШІ-генерованих текстів; етика використання ШІ в академічному письмі; плагіат та ШІ & Інформаційно-аналітичний, Рефлексивний \\
\addlinespace
Філософія & Історія філософії; онтологія; епістемологія & Базовий: модуль <<Філософія штучного інтелекту>>~--- свідомість, авторство, творчість; ШІ та творчість; відповідальність за ШІ-рішення; тест Тюрінга та creative Turing test & Рефлексивний \\
\addlinespace
Дизайн та підприємництво & Бізнес-планування; маркетинг для дизайнерів & Базовий: ШІ-інструменти для маркетингу (ChatGPT для копірайтингу, Midjourney для контенту); ШІ як конкурентна перевага; бізнес-моделі ШІ-дизайну & Комунікативний, Інформаційно-аналітичний \\
\bottomrule
\end{tabularx}
\end{table}

Особливо важливою є трансформація курсу \textit{Іноземна мова}. Аналіз опитувань роботодавців (п.~\ref{subsec:2.2}) виявив, що знання іноземної мови має найнижчу оцінку (3,0--3,5 із 5), що обмежує доступ випускників до міжнародних ШІ-ресурсів та документації. Оскільки переважна більшість ШІ-інструментів працює англійською мовою, а промпти англійською дають кращі результати, формування англомовної компетентності є необхідною передумовою ефективного формування технологічного компоненту ІКК.

\subsubsection{Нові ШІ-орієнтовані дисципліни}\label{sssec:3.2.8}

Окрім оновлення існуючих дисциплін, модель передбачає введення нових курсів, специфічно спрямованих на формування ІКК. Ці курси відповідають міжнародному бенчмаркінгу (п.~\ref{subsec:2.3}): всі 7 проаналізованих міжнародних програм мають від 1 до 4 спеціалізованих ШІ-дисциплін.

\begin{table}[htbp]
\centering
\caption{Нові ШІ-орієнтовані дисципліни}\label{tab:new-ai-courses}
\small
\begin{tabularx}{\textwidth}{@{}p{0.20\textwidth}p{0.06\textwidth}p{0.08\textwidth}p{0.27\textwidth}p{0.28\textwidth}@{}}
\toprule
\textbf{Дисципліна} & \textbf{Кр.} & \textbf{Сем.} & \textbf{Зміст} & \textbf{Компоненти ІКК} \\
\midrule
Цифрова грамотність та основи AI & 4 & 1 & Основи ШІ; типи моделей; базова промпт-інженерія; етичні основи; огляд ШІ-інструментів для дизайну & Технол.: базові навички; Інф.-ан.: розуміння ШІ; Рефлекс.: етичні основи \\
\addlinespace
Генеративний дизайн та AI & 6 & 5 & Поглиблена промпт-інженерія; Stable Diffusion (ControlNet, LoRA); параметричний дизайн; ШІ-робочі процеси; creative coding & Технол.: продвинуті навички; Інф.-ан.: оцінка; Комунік.: промптинг \\
\addlinespace
Взаємодія людина--AI в дизайні & 4 & 5 & Теорія HCI для ШІ; conversational UI; ШІ-партнерство; колаборативна творчість; роль дизайнера в ШІ-еру & Рефлекс.: глибока рефлексія; Комунік.: HCI; Інф.-ан.: аналіз взаємодії \\
\addlinespace
Критичний дизайн та етика AI & 4 & 7 & Bias в ШІ; авторське право; екологічний вплив; deepfake; speculative design; ШІ та суспільство & Рефлекс.: етична рефлексія; Інф.-ан.: критичний аналіз \\
\bottomrule
\end{tabularx}
\end{table}

Ці 4 нові дисципліни утворюють <<ШІ-спіну>> (AI spine) навчального плану~--- наскрізну послідовність від базової ШІ-грамотності (1 курс) через інструментальне використання та теоретичне осмислення (3 курс) до критичного та етичного аналізу (4 курс). Це відповідає принципу прогресивної інтеграції, визначеному як найкраща практика за результатами міжнародного бенчмаркінгу~\cite{Musiienko2025DesignEducation}.

\subsubsection{Інтегрована модель 240-кредитного навчального плану}\label{sssec:3.2.9}

На основі аналізу прогалин (розділ~\ref{sec:2}) та рекомендацій щодо трансформації дисциплін розроблено модель 4-річного (240~кредитів ЄКТС) навчального плану з наскрізною ШІ-інтеграцією. Модель адаптовано з урахуванням специфіки КДПУ та вимог стандарту В2 Дизайн.

\textbf{Рік 1. Фундамент: традиційні навички + цифрова грамотність + ШІ-ознайомлення} (60~кредитів). Фокус першого року~--- формування базових дизайнерських навичок та ШІ-грамотності. ШІ-інтеграція мінімальна (базовий рівень для більшості дисциплін), за винятком спеціалізованої дисципліни <<Цифрова грамотність та основи AI>>.

\begin{longtable}{@{}p{0.30\textwidth}p{0.06\textwidth}p{0.06\textwidth}p{0.08\textwidth}p{0.40\textwidth}@{}}
\caption{Навчальний план. Рік~1 (60~кредитів)}\label{tab:year1} \\
\hline
\textbf{Дисципліна} & \textbf{Кр.} & \textbf{Сем.} & \textbf{Рівень ШІ} & \textbf{ШІ-компонент} \\
\hline
\endfirsthead
\hline
\textbf{Дисципліна} & \textbf{Кр.} & \textbf{Сем.} & \textbf{Рівень ШІ} & \textbf{ШІ-компонент} \\
\hline
\endhead
Основи візуальної комунікації & 6 & 1 & --- & Традиційний курс \\
Рисунок і живопис & 6 & 1 & Баз. & ШІ-референси; порівняння <<ручне vs. ШІ>> \\
Історія мистецтва та дизайну & 4 & 1 & Баз. & ШІ <<у стилі>>; критичний аналіз \\
\textbf{Цифрова грамотність та основи AI} & \textbf{4} & \textbf{1} & \textbf{Транс.} & \textbf{Основи ШІ; промпт-інженерія; етика} \\
Типографіка & 5 & 1 & Баз. & ШІ-генерація шрифтових пар \\
Іноземна мова & 3 & 1 & Баз. & ШІ-термінологія англійською \\
\addlinespace
Колористика & 5 & 2 & Баз. & ШІ-палітри; кольоровий bias \\
Комп'ютерна графіка (Adobe Suite) & 6 & 2 & Інстр. & Adobe Firefly; Generative Fill \\
Дизайн-мислення & 4 & 2 & Баз. & ШІ на етапі ідеації \\
Основи композиції & 5 & 2 & Інстр. & ШІ-генерація варіантів композиції \\
Академічне письмо & 3 & 2 & Баз. & ШІ як асистент; етика \\
Вибіркова 1 & 4 & 2 & --- & --- \\
\addlinespace
\multicolumn{5}{@{}l}{\textit{Підсумок: 7~дисциплін~--- базовий рівень ШІ, 2~--- інструментальний, 1~--- трансформаційний}} \\
\hline
\end{longtable}

\textbf{Рік 2. Фахові навички: цифрові інструменти + UX/UI + моушн-дизайн} (60~кредитів). На другому курсі акцент зміщується на опанування професійних цифрових інструментів з інтеграцією ШІ-функцій. Рівень ШІ-інтеграції зростає до інструментального для більшості фахових дисциплін.

\begin{longtable}{@{}p{0.30\textwidth}p{0.06\textwidth}p{0.06\textwidth}p{0.08\textwidth}p{0.40\textwidth}@{}}
\caption{Навчальний план. Рік~2 (60~кредитів)}\label{tab:year2} \\
\hline
\textbf{Дисципліна} & \textbf{Кр.} & \textbf{Сем.} & \textbf{Рівень ШІ} & \textbf{ШІ-компонент} \\
\hline
\endfirsthead
\hline
\textbf{Дисципліна} & \textbf{Кр.} & \textbf{Сем.} & \textbf{Рівень ШІ} & \textbf{ШІ-компонент} \\
\hline
\endhead
Графічний дизайн & 6 & 3 & Інстр. & ШІ для логотипів; промпт-інженерія \\
UX/UI дизайн: Основи & 6 & 3 & Інстр. & Figma AI; ШІ-wireframes; UX з ШІ \\
3D-моделювання & 5 & 3 & Інстр. & ШІ-генерація 3D; ШІ-текстурування \\
Фотографія для дизайнерів & 4 & 3 & Баз. & ШІ-обробка; генеративне редагування \\
Веб-технології & 4 & 3 & Баз. & ШІ-асистенти для коду \\
Вибіркова 2 & 5 & 3 & --- & --- \\
\addlinespace
Figma та прототипування & 5 & 4 & Інстр. & Figma AI; ШІ-плагіни \\
Моушн-дизайн & 5 & 4 & Інстр. & RunwayML; ШІ-анімація \\
Ілюстрація & 5 & 4 & Інстр. & ШІ-ілюстрація; авторський стиль vs. ШІ \\
Інформаційний дизайн & 4 & 4 & Інстр. & ШІ-візуалізація даних \\
Навчальна практика & 6 & 4 & Інстр. & Фіксація ШІ на підприємстві \\
Вибіркова 3 & 5 & 4 & --- & --- \\
\addlinespace
\multicolumn{5}{@{}l}{\textit{Підсумок: 2~--- базовий, 8~--- інструментальний, 0~--- трансформаційний}} \\
\hline
\end{longtable}

\textbf{Рік 3. Поглиблення: ШІ-інтеграція + спеціалізація + виробнича практика} (60~кредитів). Третій курс є ключовим для формування ІКК: вводяться дві нові ШІ-дисципліни (<<Генеративний дизайн та AI>> та <<Взаємодія людина--AI в дизайні>>), а виробнича практика передбачає реальне використання ШІ у професійному середовищі.

\begin{longtable}{@{}p{0.30\textwidth}p{0.06\textwidth}p{0.06\textwidth}p{0.08\textwidth}p{0.40\textwidth}@{}}
\caption{Навчальний план. Рік~3 (60~кредитів)}\label{tab:year3} \\
\hline
\textbf{Дисципліна} & \textbf{Кр.} & \textbf{Сем.} & \textbf{Рівень ШІ} & \textbf{ШІ-компонент} \\
\hline
\endfirsthead
\hline
\textbf{Дисципліна} & \textbf{Кр.} & \textbf{Сем.} & \textbf{Рівень ШІ} & \textbf{ШІ-компонент} \\
\hline
\endhead
\textbf{Генеративний дизайн та AI} & \textbf{6} & \textbf{5} & \textbf{Транс.} & \textbf{Промпт-інженерія; Stable Diffusion} \\
\textbf{Взаємодія людина--AI в дизайні} & \textbf{4} & \textbf{5} & \textbf{Транс.} & \textbf{HCI для ШІ; колаборативна творчість} \\
Дизайн бренду: Продвинутий & 5 & 5 & Інстр. & ШІ у брендингу; ШІ-стратегії \\
Параметричний та інтерактивний дизайн & 5 & 5 & Інстр. & Алгоритмічний дизайн з ШІ \\
Дослідження в дизайні & 4 & 5 & Інстр. & ШІ для аналізу даних \\
Вибіркова 4 & 6 & 5 & --- & --- \\
\addlinespace
UX/UI: Продвинутий & 5 & 6 & Транс. & Повний UX/UI процес з ШІ \\
Візуалізація даних & 5 & 6 & Інстр. & ШІ-візуалізація; автоматизація \\
Дизайн середовища & 4 & 6 & Баз. & ШІ-візуалізація інтер'єрів \\
Виробнича практика & 9 & 6 & Транс. & ШІ на підприємстві; ШІ-портфоліо \\
Вибіркова 5 & 7 & 6 & --- & --- \\
\addlinespace
\multicolumn{5}{@{}l}{\textit{Підсумок: 1~--- базовий, 4~--- інструментальний, 4~--- трансформаційний}} \\
\hline
\end{longtable}

Ключовими інноваціями 3 курсу є два спеціалізовані ШІ-курси: <<Генеративний дизайн та AI>> (6~кредитів)~--- профільний курс з практичного опанування генеративних ШІ-інструментів, та <<Взаємодія людина--AI в дизайні>> (4~кредити)~--- теоретичний курс з колаборації людини та ШІ. Це відповідає міжнародному бенчмарку: провідні програми мають від 1 до 4 спеціалізованих ШІ-дисциплін (п.~\ref{subsec:2.3}).

\textbf{Рік 4. Синтез: дипломний проєкт + професійна підготовка + ШІ-етика} (60~кредитів). Четвертий курс орієнтований на синтез усіх набутих компетенцій у дипломному проєкті та підготовку до професійної діяльності. Введення курсу <<Критичний дизайн та етика AI>> завершує <<ШІ-спіну>> навчального плану.

\begin{longtable}{@{}p{0.30\textwidth}p{0.06\textwidth}p{0.06\textwidth}p{0.08\textwidth}p{0.40\textwidth}@{}}
\caption{Навчальний план. Рік~4 (60~кредитів)}\label{tab:year4} \\
\hline
\textbf{Дисципліна} & \textbf{Кр.} & \textbf{Сем.} & \textbf{Рівень ШІ} & \textbf{ШІ-компонент} \\
\hline
\endfirsthead
\hline
\textbf{Дисципліна} & \textbf{Кр.} & \textbf{Сем.} & \textbf{Рівень ШІ} & \textbf{ШІ-компонент} \\
\hline
\endhead
Дипломне проектування I & 10 & 7 & Транс. & ШІ як дослідницький/проєктний інструмент \\
Портфоліо та професійна практика & 4 & 7 & Інстр. & Цифрове портфоліо з ШІ-проєктами \\
Дизайн та підприємництво & 4 & 7 & Баз. & ШІ-інструменти для бізнесу \\
\textbf{Критичний дизайн та етика AI} & \textbf{4} & \textbf{7} & \textbf{Транс.} & \textbf{Bias; авторство; deepfake; екологія} \\
Вибіркова спеціалізації 1 & 8 & 7 & --- & --- \\
\addlinespace
Дипломне проектування II & 15 & 8 & Транс. & ШІ-інтеграція в дипломний проєкт \\
Переддипломна практика & 6 & 8 & Транс. & ШІ-стратегія для диплому \\
Вибіркова спеціалізації 2 & 9 & 8 & --- & --- \\
\addlinespace
\multicolumn{5}{@{}l}{\textit{Підсумок: 1~--- базовий, 1~--- інструментальний, 4~--- трансформаційний}} \\
\hline
\end{longtable}

На 4 курсі особливу роль відіграє дисципліна <<Критичний дизайн та етика AI>> (4~кредити), що завершує формування рефлексивного компоненту ІКК. Дипломне проектування (I~+~II = 25~кредитів) є інтегративним простором, де здобувач демонструє сформованість усіх 4 компонентів ІКК.

\subsubsection{Зведена характеристика навчального плану}\label{sssec:3.2.10}

\begin{table}[htbp]
\centering
\caption{Зведена характеристика ШІ-інтеграції у навчальному плані}\label{tab:curriculum-summary}
\small
\begin{tabularx}{\textwidth}{@{}Xrrrr@{}}
\toprule
\textbf{Показник} & \textbf{Рік 1} & \textbf{Рік 2} & \textbf{Рік 3} & \textbf{Рік 4} \\
\midrule
Загальна кількість дисциплін & 12 & 12 & 11 & 8 \\
Дисципліни з ШІ-інтеграцією & 10 & 10 & 9 & 6 \\
\addlinespace
Базовий рівень ШІ & 7 & 2 & 1 & 1 \\
Інструментальний рівень & 2 & 8 & 4 & 1 \\
Трансформаційний рівень & 1 & 0 & 4 & 4 \\
\addlinespace
Нові ШІ-дисципліни & 1 & 0 & 2 & 1 \\
Кредити ЄКТС & 60 & 60 & 60 & 60 \\
\bottomrule
\end{tabularx}
\end{table}

Загальна характеристика навчального плану:
\begin{itemize}
    \item \textbf{240~кредитів ЄКТС}, 4 роки, 8 семестрів~--- відповідає стандарту В2 Дизайн;
    \item \textbf{4 нові ШІ-дисципліни} (18~кредитів, 7,5\% від загального обсягу)~--- відповідає міжнародному бенчмарку (1--4 ШІ-дисципліни);
    \item \textbf{35 із 43 дисциплін} (81,4\%) мають визначений рівень ШІ-інтеграції~--- наскрізна інтеграція;
    \item \textbf{Прогресивне зростання} складності ШІ-інтеграції: від переважно базового (рік~1) через інструментальний (рік~2) до трансформаційного (роки~3--4);
    \item \textbf{Збереження фундаменту}: Рисунок, Живопис, Основи композиції, Колористика збережені як обов'язкові дисципліни з мінімальним ШІ-втручанням;
    \item \textbf{<<ШІ-спіна>>}: послідовність <<Цифрова грамотність та основи AI>> (сем.~1) $\to$ <<Генеративний дизайн та AI>> + <<Взаємодія людина--AI>> (сем.~5) $\to$ <<Критичний дизайн та етика AI>> (сем.~7) забезпечує наскрізну ШІ-підготовку.
\end{itemize}

\subsubsection{Розподіл формування компонентів ІКК за роками}\label{sssec:3.2.11}

Аналіз навчального плану за компонентами ІКК виявляє диференційовану динаміку формування кожного з чотирьох компонентів:

\begin{table}[htbp]
\centering
\caption{Акценти формування компонентів ІКК за роками навчання}\label{tab:ikk-by-year}
\small
\begin{tabularx}{\textwidth}{@{}Xp{0.20\textwidth}p{0.20\textwidth}p{0.20\textwidth}p{0.20\textwidth}@{}}
\toprule
\textbf{Компонент ІКК} & \textbf{Рік 1} & \textbf{Рік 2} & \textbf{Рік 3} & \textbf{Рік 4} \\
\midrule
Інформаційно-аналітичний & ШІ-грамотність; критичний аналіз & Порівняльний аналіз інструментів & ШІ-дослідження; аналіз даних & Аналіз для дипломного проєкту \\
\addlinespace
Комунікативний & Англомовна ШІ-термінологія & Командні проєкти; цифрові платформи & Промпт-інженерія як комунікація; HCI & Портфоліо; професійна презентація \\
\addlinespace
Технологічний & Базові ШІ-інструменти; Adobe Firefly & UX/UI з ШІ; 3D з ШІ; моушн з ШІ & Генеративний дизайн; Stable Diffusion & Дипломний проєкт з ШІ \\
\addlinespace
Рефлексивний & Етичні основи ШІ & Рефлексія <<ручне vs. ШІ>> & Глибока рефлексія; maker vs. curator & ШІ-етика; критичний дизайн \\
\bottomrule
\end{tabularx}
\end{table}

Таблиця~\ref{tab:ikk-by-year} демонструє, що всі 4 компоненти ІКК формуються протягом усього навчання, проте з різними акцентами:
\begin{itemize}
    \item \textbf{технологічний} компонент зростає від базового (рік~1) до продвинутого (рік~3) та синтетичного (рік~4);
    \item \textbf{рефлексивний} компонент розвивається від етичних основ (рік~1) через практичну рефлексію (роки~2--3) до глибокого критичного аналізу (рік~4);
    \item \textbf{інформаційно-аналітичний} компонент еволюціонує від базової грамотності (рік~1) до дослідницької аналітики (рік~3);
    \item \textbf{комунікативний} компонент розвивається від термінологічної бази (рік~1) через командну роботу (рік~2) та промпт-інженерію як комунікацію (рік~3) до професійного портфоліо (рік~4).
\end{itemize}

Така диференціація забезпечує \textit{прогресивне формування ІКК} відповідно до дидактичних принципів, визначених у структурно-функціональній моделі (п.~\ref{subsec:3.1}). Розроблені рекомендації потребують визначення педагогічних умов та стратегії впровадження, що є предметом наступного підрозділу.
